## Objective

The objective of this work is to investigate whether coherent control of nuclear states in fissile material can, in principle, be used to modify the effective probability or channel distribution of nuclear fission and thereby introduce a new control variable for macroscopic fission systems.

Current nuclear reactor control relies on classical macroscopic mechanisms, such as neutron absorption, moderation, geometry, temperature feedback, and fuel composition. Quantum-mechanical effects enter reactor physics primarily through nuclear data and reaction cross sections, but the quantum state of the fissile nuclei is not itself treated as an actively controllable degree of freedom.

This work explores a different conceptual framework: the possibility of preparing or perturbing nuclear states through coherent external interactions and studying whether such state manipulation can produce a measurable modification of fission observables, including effective fission probability, branching ratios, neutron yield, or fragment distributions.

The central research question is therefore:

**Can the quantum state of fissile matter be promoted from a microscopic internal variable to an externally controllable parameter capable of influencing macroscopic fission dynamics?**

The aim is not to propose an immediately realizable reactor technology, but to establish the physical conditions, theoretical limits, and possible control mechanisms under which coherent nuclear-state manipulation could become relevant to fission control.

If such a coupling exists and is sufficiently strong, it would define a control paradigm distinct from conventional reactor physics, in which reactor behavior is influenced not only through the neutron population and macroscopic material properties, but also through controlled modification of the microscopic nuclear state of the fuel.
